\section{The exact cutoff for the logarithmic-integral inequality}
\label{sec:li-cut}
MathWorld gives Ramanujan's inequality for $\pi$ and a companion inequality for the principal-value logarithmic integral. The companion has an exact integer transition. Monotonicity reduces the whole infinite range to two rational sign certificates.

\begin{theorem}
For integers $n\ge2$,
\[
 \operatorname{li}(n)^2 <
 \frac{en}{\log n}\operatorname{li}(n/e)
 \quad\Longleftrightarrow\quad n\ge2418 .
\]
The inequality is reversed strictly for $2\le n\le2417$.
\end{theorem}
\begin{proof}
For $x>e$ put $l=\log x$, $f(x)=\operatorname{li}(x)$, and
\[
 D(x)=e f(x/e)-\frac{l}{x} f(x)^2 .
\]
Multiplying $D(x)$ by $x/l>0$ gives the right side of the claimed inequality minus its left side. The derivative $f'(x)=1/l$ gives
\[
 D'(x)=\frac1{l-1}-\frac{2f(x)}x+(l-1)\frac{f(x)^2}{x^2}
       =\frac{(1-(l-1)f(x)/x)^2}{l-1}\ge0.
\]
Consequently it suffices to determine the signs at $2417$ and $2418$ on the interval $[3,\infty)$. For $n=2$, $0<2/e<1$ and $f(2/e)<0$ by integration of a strictly negative integrand. Strictness of the two endpoint certificates below suffices even if the square in the derivative vanishes somewhere.

Here are the definitions and numerical error bounds behind the sign test. With $t=\log x$ write
\[
 \operatorname{Ei}(t)=\operatorname{li}(x)
   =\operatorname{PV}\int_{-\infty}^{t} e^u\,du/u .
\]
Substitution $x=e^u$ gives the displayed normalization. Symmetric excisions at $x=1$ or at $u=0$ have the same limit because $\log(1+h)/[-\log(1-h)]\to1$; this follows by sandwiching each log between $h/(1+h)$ and $h/(1-h)$. On finite intervals the Taylor series of $(e^u-1)/u$ can be integrated termwise with the following explicit bound. Write $S_K(t)=\sum_{j=1}^{K}t^j/(j\,j!)$. For $0<t\le B=32$ and $K=256$,
\[
 \operatorname{Ei}(t)=\log(t/B)+S_K(t)-S_K(-B)+E,\qquad
 -\frac{2^{-B}}{B}-\epsilon\le E\le\epsilon,\quad
 \epsilon=\frac{4B^{K+1}}{(K+1)!}.
\]
Indeed the integral to $-B$ is negative and has size at most $e^{-B}/B$. The principal value of $1/u$ from $-B$ to $t$ is $\log(t/B)$. Each series tail after integration is bounded by $2B^{K+1}/(K+1)!$ since successive absolute terms have ratio at most $1/2$. These arguments also provide arbitrary-precision names by increasing $B$ and $K$.

To evaluate logarithms normalize a positive rational to $2^a y$ with $1\le y\le2$. With $z=(y-1)/(y+1)$ integrate the finite geometric series for $2/(1-z^2)$. This gives
\[
 0\le \log y-2\sum_{j=0}^{49}\frac{z^{2j+1}}{2j+1}
 \le \frac{2z^{101}}{101(1-z^2)}.
\]
The same rule evaluates $\log 2$. Evaluate $e$ by summing through $1/50!$ with tail at most $2/51!$. For $t=\log n$, the quantity $nD(n)=en\operatorname{Ei}(t-1)-t\operatorname{Ei}(t)^2$ has the following enclosures. All endpoints in the table are integers divided by $1000$.
\[
\begin{array}{c|rr}
 n & \text{strict lower numerator}& \text{strict upper numerator}\\\hline
 2417 & -437 &-435\\
 2418 & 40 &42
\end{array}
\]
The file \texttt{certify\_li.py} reproduces these rational comparisons. It accumulates $t^j/j!$ by multiplication by $t/j$ and rounds the two endpoints out to multiples of $2^{-120}$. It applies the logarithm bounds at each end and uses monotonicity of $\operatorname{Ei}$ on the positive axis. Thus each line is a closed finite calculation, reproducible using integer arithmetic.

For clarity, the monotonicity step also has a finite interpretation. Fix integers $3\le a\le b$. On $[3,b]$, $\log x-1\ge1/16$; already the positive terms $2(1/2+(1/2)^3/3)$ of the logarithm series for $3$ exceed $1+1/16$. We have $|\operatorname{li}(x)|\le b+10$ there. Indeed the same expansion at $t=\log3$ gives $|\operatorname{li}(3)|<10$ (also checked by the auxiliary script), and $|f'|\le1$ on $[3,b]$. Write $A=1-(l-1)f/x$. The elementary derivative bounds give $|A|\le 1+b(b+10)$ and $|A'|\le (b+1)(b+11)$ using $x\ge3$, $l\le b$. Differentiating $A^2/(l-1)$ thus gives $|D''|\le M=10^4(b+12)^5$. On a mesh of width $h$ the error of the left sum of $D'$ for $D(b)-D(a)$ is at most $M(b-a)h$. This estimate can also be derived without a mean-value witness: differences of the integral names $f$ and $\log$ on $[u,u+h]$ equal their left derivatives times $h$ with quadratic remainders, by upper and lower rectangles for their decreasing positive kernels. Applying the product and reciprocal difference rules gives the same first-order expansion for $D$; equivalently integrate the bounded derivative of $D'$ cell by cell. Take $h\le2^{-s-2}/(M(b+1))$ and give each summand an accumulated rounding allowance below $2^{-s-2}$. The sum of squares is nonnegative, so $D(b)\ge D(a)-2^{-s}$ at every requested precision. Fixed positive margins in the two endpoint tests turn this into the stated strict inequalities. The mesh, all apartness bounds, and the power-series lengths are primitive recursive in the endpoints and precision.
\end{proof}

For an explicit general series schedule in the preceding paragraph, let the required absolute error be $2^{-s}$, take $B\ge s+\lceil|t|\rceil+8$, and use $K+1=2m$ with $m=4(B+1)^2+s+8$. Since $(2m)!\ge m^m$ and $m\ge4B^2$, the tail allowance $4 B^{2m}/(2m)!$ is at most $4^{1-m}$. The logarithm at $t>0$ in these applications is evaluated with the known lower bound on $t$ (at $x/e$ it is at least $1/16$ for $x\ge3$). General dyadic widths can be propagated by the finite circuit estimate of Section~\ref{sec:mellin}.
